\documentclass[12pt,a4paper]{article}

\usepackage[a4paper,margin=1in]{geometry}
\usepackage{graphicx}
\usepackage{booktabs}
\usepackage{amsmath,amssymb}
\usepackage{float}
\usepackage{hyperref}
\usepackage{xcolor}
\usepackage{listings}
\usepackage{enumitem}
\usepackage{fancyhdr}
\usepackage{longtable}

\hypersetup{colorlinks=true,linkcolor=blue,urlcolor=blue}

\cfoot{\thepage}

\lstset{
language=Python,
basicstyle=\ttfamily\small,
numbers=left,
frame=single,
breaklines=true,
showstringspaces=false
}

\begin{document}

\begin{tabular}{|p{4cm}|p{11cm}|}
\hline
Experiment No. & 6 \\ \hline
Experiment Title & Dimensionality Reduction and Model Evaluation (With and Without PCA)\\ \hline
Student Name & S.S.Janane\\ \hline
Register Number & 3122247001024 \\ \hline
Faculty & Dr. Poreddy Ajay Kumar Reddy \\ \hline
Submission Date & 30-8-2026\\ \hline
\end{tabular}

\section{Objective}
This experiment studies how Principal Component Analysis (PCA) affects the performance of ten different classifiers. Each model is trained and validated twice, once on the original features and once on the PCA-reduced features, using hyperparameter tuning and 5-fold cross-validation in both cases, so that the two settings can be compared fairly.

\section{Problem Statement}
Given a set of numerical features that describe breast tissue samples, the task is to classify each sample as malignant or benign. The input is a set of 30 numerical features per ,sample and the output is a binary label. The goal is to check whether reducing the number of features through PCA changes accuracy, F1-score, or stability across different types of models.

\section{Dataset Description}
\begin{longtable}{|p{6cm}|p{8cm}|}
\hline
Dataset Name & Wisconsin Diagnostic Breast Cancer (WDBC) \\ \hline
Dataset Source & UCI Machine Learning Repository, accessed through scikit-learn's load\_breast\_cancer \\ \hline
Number of Samples & 569 \\ \hline
Number of Features & 30 numerical features \\ \hline
Number of Classes & 2 (Malignant, Benign) \\ \hline
Class Distribution & 357 Benign, 212 Malignant \\ \hline
Missing Values & None \\ \hline
Train-Test Split & 80\% train (455 samples), 20\% test (114 samples), stratified \\ \hline
\end{longtable}

\begin{figure}[H]
\centering
\includegraphics[width=0.55\linewidth]{class_distribution.png}
\caption{Class distribution of Malignant vs Benign samples}
\end{figure}

\subsection*{Inference}
The data set is slightly imbalanced, with more benign than malignant samples. Stratified splitting was used for the train-test split and for every cross-validation fold, so this ratio is preserved throughout the experiment and does not bias the comparison between models.

\section{Data Preprocessing}
\begin{itemize}
\item Missing value handling: not required the dataset has no missing values
\item Encoding: not required all features are numerical and the target is binary
\item Feature scaling: StandardScaler was fit to the training data only and applied to both the train and the test sets
\item Train-test split: 80-20 split with stratification on the target class
\item PCA: applied on top of the scaled features, described in the next section
\end{itemize}

\section{Principal Component Analysis}

\begin{longtable}{|p{4cm}|p{5cm}|p{2.5cm}|p{4cm}|}
\hline
Setting & Chosen Components / Variance Target & Explained Variance (\%) & Justification \\ \hline
With-PCA & 10 components (95\% variance target) & 95.27 & Keeps most of the information while roughly cutting the feature count in half \\ \hline
\end{longtable}

\begin{figure}[H]
\centering
\includegraphics[width=\linewidth]{scree_plot.png}
\caption{Scree plot: individual and cumulative explained variance by component}
\end{figure}

\subsection*{Inference}
The individual variance plot shows that the first few components carry most of the information, and the curve flattens out quickly after that. The cumulative plot shows that 10 of the original 30 components are enough to cross the 95\% variance mark, so the remaining 20 components add very little additional information.

\section{Machine Learning Algorithms}
Ten classifiers were evaluated: Support Vector Machine, Naive Bayes, k-Nearest Neighbors, Logistic Regression, Decision Tree, Random Forest, AdaBoost, Gradient Boosting, XGBoost and a Stacked Ensemble of SVM, Random Forest, and KNN with a Logistic Regression meta-learner. SVM and Logistic Regression are linear or margin-based models that work directly with the geometry of the feature space, so they are expected to respond differently to PCA compared to tree-based models such as Decision Tree, Random Forest, and the boosting methods, which build their decision boundaries by splitting on individual features rather than relying on distances or linear combinations.

\section{Experimental Setup}
\begin{tabular}{ll}
\toprule
Python Version & 3.11 \\
Scikit-Learn Version & 1.x \\
IDE & Jupyter Notebook \\
Operating System & Linux \\
Hardware & Standard CPU, no GPU used \\
\bottomrule
\end{tabular}

\section{Hyperparameter Selection}

\begin{longtable}{lll}
\toprule
Model & Best Params (No-PCA) & Best Params (With-PCA) \\
\midrule
\endhead
SVM & C=0.1, gamma=scale, kernel=linear & C=10, gamma=0.01, kernel=rbf \\
Naive Bayes & var\_smoothing=1e-9 & var\_smoothing=1e-9 \\
KNN & k=3, manhattan, distance & k=5, euclidean, distance \\
Logistic Regression & C=0.1 & C=0.1 \\
Decision Tree & max\_depth=3, min\_samples\_split=5 & max\_depth=5, min\_samples\_split=2 \\
Random Forest & max\_depth=10, n\_estimators=100 & max\_depth=None, n\_estimators=200 \\
AdaBoost & learning\_rate=1.0, n\_estimators=100 & learning\_rate=1.0, n\_estimators=50 \\
Gradient Boosting & learning\_rate=0.1, max\_depth=2, n=100 & learning\_rate=0.1, max\_depth=3, n=200 \\
XGBoost & learning\_rate=0.05, max\_depth=3, n=100 & learning\_rate=0.05, max\_depth=5, n=200 \\
Stacking & final\_estimator C=1 & final\_estimator C=1 \\
\bottomrule
\end{longtable}

All hyperparameters were selected using 5-fold stratified cross-validation on the training set, separately for the No-PCA and With-PCA feature sets.

\section{Results}

\subsection*{5-Fold Cross-Validation Results}

\begin{longtable}{lccccccc}
\toprule
Model & Fold 1 & Fold 2 & Fold 3 & Fold 4 & Fold 5 & Avg No-PCA & Avg With-PCA \\
\midrule
\endhead
SVM & 95.60 & 97.80 & 97.80 & 97.80 & 98.90 & 97.58 & 98.02 \\
Naive Bayes & 91.21 & 96.70 & 89.01 & 95.60 & 94.51 & 93.41 & 91.65 \\
KNN & 97.80 & 98.90 & 95.60 & 95.60 & 97.80 & 97.14 & 96.48 \\
Logistic Regression & 97.80 & 97.80 & 98.90 & 97.80 & 98.90 & 98.24 & 98.24 \\
Decision Tree & 92.31 & 93.41 & 92.31 & 92.31 & 93.41 & 92.75 & 92.31 \\
Random Forest & 96.70 & 95.60 & 93.41 & 96.70 & 98.90 & 96.26 & 96.04 \\
AdaBoost & 98.90 & 100.00 & 94.51 & 97.80 & 98.90 & 98.02 & 95.16 \\
Gradient Boosting & 95.60 & 98.90 & 94.51 & 96.70 & 97.80 & 96.70 & 96.26 \\
XGBoost & 98.90 & 96.70 & 94.51 & 96.70 & 98.90 & 97.14 & 96.48 \\
Stacking & 97.80 & 96.70 & 95.60 & 96.70 & 98.90 & 97.14 & 97.14 \\
\bottomrule
\end{longtable}

The fold values shown above are for the No-PCA setting, and the last two columns give the average accuracy for No-PCA and With-PCA respectively, so the change with PCA can be read directly from those two columns.

\subsection*{Test Set Results (Selected Models)}

\begin{tabular}{llccccc}
\toprule
Model & Setting & Accuracy (\%) & Precision & Recall & F1 & AUC \\
\midrule
SVM & No-PCA & 98.25 & 0.986 & 0.986 & 0.986 & 0.994 \\
SVM & With-PCA & 95.61 & 0.986 & 0.944 & 0.965 & 0.996 \\
Random Forest & No-PCA & 95.61 & 0.959 & 0.972 & 0.966 & 0.994 \\
Random Forest & With-PCA & 93.86 & 0.958 & 0.944 & 0.951 & 0.986 \\
XGBoost & No-PCA & 94.74 & 0.946 & 0.972 & 0.959 & 0.993 \\
XGBoost & With-PCA & 93.86 & 0.958 & 0.944 & 0.951 & 0.993 \\
Stacking & No-PCA & 97.37 & 0.986 & 0.972 & 0.979 & 0.994 \\
Stacking & With-PCA & 95.61 & 0.972 & 0.958 & 0.965 & 0.991 \\
\bottomrule
\end{tabular}

\section{Result Visualization}

\subsection*{Cross-Validation Accuracy Comparison}

\begin{figure}[H]
\centering
\includegraphics[width=\linewidth]{cv_accuracy_comparison.png}
\caption{Average 5-fold CV accuracy for each model, No-PCA vs With-PCA}
\end{figure}

\subsection*{Inference}
Most models sit close together, whether PCA is applied or not. SVM is the only model where the With-PCA bar is clearly taller than the No-PCA bar. AdaBoost shows the biggest drop, with its accuracy falling by close to three percentage points once PCA is applied.

\subsection*{Fold-to-Fold Variability}

\begin{figure}[H]
\centering
\includegraphics[width=\linewidth]{fold_variability.png}
\caption{Standard deviation of fold accuracy for each model, No-PCA vs With-PCA}
\end{figure}

\subsection*{Inference}
For several models, including Decision Tree, Gradient Boosting, and AdaBoost, the spread across folds becomes larger with PCA rather than smaller. Only Logistic Regression and Stacking keep exactly the same spread in both settings, since their best hyperparameters and scores did not change.

\subsection*{Confusion Matrices}

\begin{figure}[H]
\centering
\includegraphics[width=\linewidth]{confusion_matrices.png}
\caption{Confusion matrices for SVM, Random Forest, XGBoost and Stacking, No-PCA vs With-PCA}
\end{figure}

\subsection*{Inference}
For all four models shown, the No-PCA version makes fewer errors or an equal number compared to the With-PCA version on the test set. SVM without PCA gets the closest to a perfect confusion matrix, while the With-PCA versions show a few more malignant cases that are misclassified as benign.

\subsection*{ROC Curves}

\begin{figure}[H]
\centering
\includegraphics[width=0.7\linewidth]{roc_curve.png}
\caption{ROC curves for SVM, Random Forest, XGBoost and Stacking, No-PCA vs With-PCA}
\end{figure}

\subsection*{Inference}
All eight curves stay close to the top left corner, which means that every model ranks the two classes well regardless of PCA. The AUC values are all above 0.98, and the small differences between the No-PCA and With-PCA curves for the same model are much smaller than the differences between different model types.

\section{Discussion}
Of the ten models tested, PCA only gave a clear benefit to SVM, whose cross-validation accuracy increased from 97.58\% to 98.02\%. Logistic Regression and Stacking were not affected, staying at exactly the same accuracy in both settings. The rest of the models, especially AdaBoost and Naive Bayes, performed worse with PCA. PCA also did not make the results more stable, the fold-to-fold spread went up for more models than it went down for. This line up with the fact that WDBC has only 30 features to begin with, which is not extremely high-dimensional, so there was limited redundant information for PCA to usefully remove for most model types. Margin-based models such as SVM benefitted because PCA components are uncorrelated, which helps the margin computation, while boosting and probability-based models lost some of the specific feature signal on which they were relying.

\section{Conclusion}
Ten classifiers were tuned and evaluated with 5-fold cross-validation in the WDBC data set, both without PCA and with PCA reduced to 10 components covering 95.27\% of the variance. PCA improved only SVM, left Logistic Regression and Stacking unchanged, and reduced accuracy for the remaining seven models, most noticeably AdaBoost. The stacked ensemble was the most resistant to the change, as its accuracy did not change at all despite its individual base learners responding differently. Overall, PCA is not a universally helpful preprocessing step here, its benefit depends heavily on the type of model being used.

\section{References}
\begin{enumerate}
\item I. T. Jolliffe, ``Principal Component Analysis,'' Springer Series in Statistics, 2002.
\item L. Breiman, ``Random Forests,'' Machine Learning, vol. 45, no. 1, pp. 5--32, 2001.
\item T. Chen and C. Guestrin, ``XGBoost: A Scalable Tree Boosting System,'' Proc. 22nd ACM SIGKDD, pp. 785--794, 2016.
\item Scikit-learn Developers, ``Decomposition: PCA,'' [Online]. Available: \url{https://scikit-learn.org/stable/modules/decomposition.html#pca}
\end{enumerate}

\end{document}
